\section{Introduction}\label{sec:intro}

Modern AI agents spend much of their time on decisions that are small but frequent. Should the agent
read \code{report-final.md} or \code{report-draft.md}? Is the \emph{Save} button the one the task needs?
Does this function really enforce the size limit it claims to? A large language model can answer each
of these questions well, but it takes seconds and costs money every time. The agent bundle studied here,
\emph{Fast Decisions}, hands such questions to a small, fast \term{decision judge} instead. The judge
returns a probability for each option. If it is confident enough, the agent acts at once. If not, the
agent falls back to its ordinary, slower reasoning.

This design only helps if the judge is right when it is confident. A judge that is confidently wrong
does not merely waste time: it makes the agent do the wrong thing with no second look. This report asks
a practical question: \emph{which fast model should make these decisions, and under what safeguards?}

\paragraph{What we did.} We compared \NumBase\ judges: two hosted OpenAI language models, one hosted
decision model (Jev~1.13, the bundle's current default) and \NumLocal\ models that run on the same
computer as the agent. We also tested \NumClauseArms\ judge variants with one extra sentence in their
instructions. Every judge answered the same structured questions, \DevReps\ times each, and every
answer was scored under the exact rule the bundle uses to decide whether to act.

The work had two parts.
\begin{enumerate}
  \item \textbf{Validate a first pass.} An earlier, quick comparison on the same \DevN\ development
    cases (the \term{first pass}) made \ChangesN\ claims. We reproduced it, audited its harness and
    labels, and checked each claim (\cref{sec:validate}).
  \item \textbf{Confirm on fresh data.} We then wrote \HoldN\ new cases, had their labels checked blind,
    and committed the analysis plan and decision rules \emph{before} any judge saw them
    (a \term{preregistration}). Only this \term{holdout} can confirm a claim (\cref{sec:holdout}).
\end{enumerate}

\begin{keyidea}
Accuracy alone is the wrong headline for a judge. What matters most is how often the agent
\emph{acts automatically on a wrong answer} (the wrong-automatic rate), and how often it gets to act
automatically at all (coverage). A wrong answer that falls back to slow reasoning costs only time.
\end{keyidea}

\paragraph{Main findings.} In short (details and uncertainty in \cref{sec:recs}):
\begin{itemize}
  \item \textbf{Keep Jev~1.13 as the default.} It was right on \JevHoldAccK/\HoldN\ holdout cases with
    \JevHoldWaK\ wrong automatic decision, a p95 latency of \JevHoldPninetyfive\,ms and
    \$\JevHoldCost\ per million decisions. The two OpenAI models scored higher, but those differences
    are within noise after correcting for multiple comparisons, and both are much slower
    (p95 \LunaHoldPninetyfiveS\,s and \SolHoldPninetyfiveS\,s).
  \item \textbf{Add a host guard on side-effect actions.} A fixed word list that blocks automatic
    purchases, deletions, publications and similar actions removed every side-effect wrong automatic
    decision it targeted on the holdout, in all \ITwoHoldNApplicable\ judges it applied to.
  \item \textbf{No local model is ready to be the offline tier.} The best, \OffBestHold, still made
    \TevFourHoldOffWaK\ wrong automatic decisions (upper 95\,\% bound \TevFourHoldOffWaUpper\,\%,
    above the preregistered \OffMaxUpperPct\,\% limit).
  \item \textbf{The first pass was partly wrong.} Its headline that GPT-6 Luna was right on all
    \DevN\ cases with no wrong automatic decisions did not hold up (\cref{sec:validate}).
\end{itemize}

\paragraph{How to read this report.} \Cref{sec:background} explains judges and automatic decisions
with one worked example. \Cref{sec:method} describes the cases, labels, judges and statistics in plain
terms. \Cref{sec:validate,sec:dev,sec:holdout} give the results; results from the development split
are marked \screen\ and results from the holdout are marked \prereg. \Cref{sec:failures,sec:latency}
explain \emph{why} judges fail and where their time goes. \Cref{sec:recs} gives recommendations, and
\cref{sec:limits,sec:repro} state the limits and how to reproduce everything. A glossary is in
\cref{app:glossary}.
